%! AP Statistics — Unit 2, Lesson 2.2 — Guided Notes KEY
\documentclass[10pt]{article}
\usepackage{apstats-article}
\usepackage{apstats-key}

%=============================================================================
\begin{document}

\pageheader{Unit 2, Lesson 2.2}{Guided Notes --- Answer Key}
\begin{tabularx}{\linewidth}{X X X}
Name:\;\ans{Key} & Date:\; & Period:\;
\end{tabularx}

\vspace{0.12in}

\begin{objectivebox}
\small By the end of this lesson, I will be able to\dots\\[4pt]
\ans{construct a two-way (contingency) table and identify joint and marginal frequencies;}\\[1pt]
\ans{compute joint and conditional relative frequencies from a two-way table and explain the difference;}\\[1pt]
\ans{construct and interpret a segmented bar chart to determine whether two categorical variables are associated.}
\end{objectivebox}

\vspace{0.1in}

\begin{vocabbox}
\termblanklong{Two-way table (contingency table)}\hfill\ans{a table that organizes counts for two categorical variables simultaneously; rows = one variable, columns = the other}\\
\termblanklong{Joint frequency}\hfill\ans{the count in a single interior cell; describes both variables at once}\\
\termblanklong{Marginal frequency}\hfill\ans{a row or column total; describes one variable ignoring the other}\\
\termblanklong{Joint relative frequency}\hfill\ans{a cell count divided by the grand total of the entire table (UNC-1.P.4)}\\
\termblanklong{Conditional relative frequency}\hfill\ans{a cell count divided by a row or column total; holds one variable fixed}\\
\termblanklong{Association (two categorical variables)}\hfill\ans{present when conditional distributions differ across groups; if they are the same, there is no evidence of association}
\end{vocabbox}

\vspace{0.1in}

\begin{hookbox}
\small
\begin{enumerate}[leftmargin=1.5em, itemsep=2pt]
  \item \ans{A single bar chart shows the distribution of one variable. To compare
  lunch preferences across grade levels we need to show both variables at once --- a
  two-way table does that.}

  \item \ans{A two-way (contingency) table, or a segmented bar chart built from it.}
\end{enumerate}
\end{hookbox}

\vspace{0.1in}

\begin{notesbox}{Step 1 --- Build the Two-Way Table (UNC-1.P.3)}
\small

\vspace{8pt}
\renewcommand{\arraystretch}{2.0}
\begin{tabularx}{\linewidth}{@{} l C{2.4cm} C{2.4cm} C{2.2cm} C{2.2cm} @{}}
\rowcolor{sky}
\textbf{Grade Level}
  & \textbf{Cafeteria} & \textbf{Off-campus} & \textbf{Home}
  & \textbf{\color{navy}Row Total} \\
\hline
Underclassman & 54 & 18 & 48 & \ans{120} \\
Upperclassman & 36 & 42 & 42 & \ans{120} \\
\hline
\textbf{Column Total} & \ans{90} & \ans{60} & \ans{90} & \ans{240} \\
\end{tabularx}

\vspace{10pt}

\begin{enumerate}[leftmargin=1.4em, itemsep=4pt]
  \item Row totals: 120 and 120; column totals: 90, 60, 90.
        These are called \ans{marginal} frequencies.

  \item The count in any interior cell is called a \ans{joint} frequency because
        it describes \ans{both variables simultaneously (e.g., students who are
        underclassmen AND prefer the cafeteria)}.

  \item Grand total = \ans{240}. Verify: $\ans{120} + \ans{120} = \ans{240}$ \checkmark
\end{enumerate}

\end{notesbox}

\vspace{0.1in}

\begin{notesbox}{Step 2 --- Relative Frequencies (UNC-1.P.4)}
\small
\begin{multicols}{2}

\textbf{Joint relative frequency}\\[3pt]
Formula: $\dfrac{\text{cell count}}{\ans{\text{grand total}}}$

\vspace{6pt}
\[
  \frac{\ans{54}}{\ans{240}} = \ans{0.225} \approx 22.5\%
\]

\columnbreak

\textbf{Conditional relative frequency}\\[3pt]
Formula: $\dfrac{\text{cell count}}{\ans{\text{row (or column) total}}}$

\vspace{6pt}
Among underclassmen:
\[
  \frac{\ans{54}}{\ans{120}} = \ans{0.45 = 45\%}
\]

\end{multicols}

\vspace{6pt}
\textbf{Key distinction:} Joint uses the \ans{grand} total;
conditional uses a \ans{row} or \ans{column} total.
\end{notesbox}

\vspace{0.1in}

\begin{notesbox}{Step 3 --- Conditional Distribution Table}
\small

\vspace{8pt}
\renewcommand{\arraystretch}{2.0}
\begin{tabularx}{\linewidth}{@{} l C{2.8cm} C{2.8cm} C{2.8cm} C{1.6cm} @{}}
\rowcolor{sky}
\textbf{Grade Level}
  & \textbf{Cafeteria} & \textbf{Off-campus} & \textbf{Home}
  & \textbf{Total} \\
\hline
Underclassman & \ans{45\%} & \ans{15\%} & \ans{40\%} & 100\% \\
Upperclassman & \ans{30\%} & \ans{35\%} & \ans{35\%} & 100\% \\
\end{tabularx}

\vspace{10pt}
\ans{The two rows look different: underclassmen are much more likely to eat in the
cafeteria (45\% vs.\ 30\%) and much less likely to go off-campus (15\% vs.\ 35\%).
This suggests there \textbf{is} evidence of an association between grade level and
lunch preference.}
\end{notesbox}

\vspace{0.1in}

\begin{notesbox}{Step 4 --- Segmented Bar Chart (UNC-1.P.1, UNC-1.P.2)}
\small
\begin{multicols}{2}

\begin{enumerate}[leftmargin=1.4em, itemsep=3pt]
  \item Each bar reaches \ans{100}\%.
  \item Divide each bar by \ans{conditional} relative frequencies.
  \item Label each segment (Cafeteria / Off-campus / Home).
\end{enumerate}

\vspace{6pt}
\textbf{Reading the chart:}\\
If the two bars look \textbf{the same} $\to$ \ans{no association}\\
If the two bars look \textbf{different} $\to$ \ans{evidence of association}

\columnbreak

\vspace{4pt}
\textbf{Segmented bar chart (key sketch):}

\vspace{0.5in}

\begin{tikzpicture}[x=1.4cm, y=0.035cm]
  \draw[thick,->] (0,0) -- (2.8,0);
  \draw[thick,->] (0,0) -- (0,115) node[above,font=\tiny] {\%};
  \foreach \y in {0,25,50,75,100}{
    \draw (0,\y) -- (-0.08,\y) node[left,font=\tiny] {\y};
  }
  \node[below,font=\small] at (0.7,0) {Under};
  \node[below,font=\small] at (2.1,0) {Upper};
  % Underclassman bar: Cafeteria 0-45, Off-campus 45-60, Home 60-100
  \fill[navy!70]   (0.2,  0) rectangle (1.2, 45);
  \fill[sky!80!navy] (0.2, 45) rectangle (1.2, 60);
  \fill[greenbg!80] (0.2, 60) rectangle (1.2,100);
  \draw (0.2,0) rectangle (1.2,100);
  % Upperclassman bar: Cafeteria 0-30, Off-campus 30-65, Home 65-100
  \fill[navy!70]   (1.6,  0) rectangle (2.6, 30);
  \fill[sky!80!navy] (1.6, 30) rectangle (2.6, 65);
  \fill[greenbg!80] (1.6, 65) rectangle (2.6,100);
  \draw (1.6,0) rectangle (2.6,100);
  % legend
  \fill[navy!70]   (0.0,108) rectangle (0.18,116); \node[right,font=\tiny] at (0.2,112) {Cafe};
  \fill[sky!80!navy] (0.55,108) rectangle (0.73,116); \node[right,font=\tiny] at (0.75,112) {Off};
  \fill[greenbg!80] (1.0,108) rectangle (1.18,116); \node[right,font=\tiny] at (1.2,112) {Home};
\end{tikzpicture}

\end{multicols}

\end{notesbox}

\vspace{0.1in}

\begin{notesbox}{Step 5 --- State the Association Conclusion}
\small

\ans{There \textbf{is} evidence of an association between grade level and lunch
preference. Among underclassmen, 45\% prefer the cafeteria, compared to 30\%
among upperclassmen. The conditional distributions \textbf{differ}, suggesting
that upperclassmen are more likely to eat off-campus and less likely to use the
cafeteria than underclassmen.}

\end{notesbox}

\vspace{0.1in}

\begin{spiralbox}
\small
\begin{multicols}{2}

\textbf{How this connects:}
\begin{itemize}[itemsep=3pt]
  \item Lesson 2.1: association question $\to$ today: displays
  \item Lesson 2.3: quantify with statistics
  \item Unit 8: chi-square significance test
\end{itemize}

\columnbreak

\ans{We use conditional relative frequencies because we want to compare
lunch preferences \emph{within} each grade level. Joint relative frequencies
include differences in group size, which can mislead --- for example, both
grade levels have the same number of students here, but in general they might not.
Conditional frequencies put each group on the same ``100\%'' scale.}

\end{multicols}
\end{spiralbox}

\begin{teachernote}
\textbf{Common errors:}
\begin{itemize}[itemsep=3pt]
  \item Dividing by the grand total (240) to get conditional frequencies --- redirect: ``that gives you the joint relative frequency, not the conditional.''
  \item Confusing rows and columns --- ask: ``which variable is explanatory? Condition on that one.''
  \item In the bar chart, bars that do not reach exactly 100\% --- rounding issues; accept $\pm$1\%.
  \item Stating ``there is an association'' without citing specific values --- require the comparison of conditional percentages as evidence.
\end{itemize}
\end{teachernote}

\end{document}
